\section{Fractal norms, compact arrays, and norm transport}
\label{d:section}

From this point the spatial dimension is $n\ge3$.  This section supplies
a weighted fractal estimate with an occupancy gain and a power saving
for frequency supports near affine $(n-2)$-planes.  In
Section~\ref{t:section}, these bounds will locate densely occupied
boxes and then permit pruning of the incident tube families so that
their direction laws do not concentrate near such planes.  We use
\begin{equation}
 \beta=\frac2{n+1},\qquad e=\frac{n^2}{n+1},
 \qquad \gamma_n=\frac{2n}{(n+1)(n+2)}.
 \label{d:exponents}
\end{equation}
The exponent $\beta$ is the measure homogeneity in the fractal square
norm.  Keeping that homogeneity explicit will distinguish a genuine
sparsity gain from a change of normalization.

\subsection{Arrays and measures}

\begin{definition}[Normalized compact arrays and admissible measures]
\label{d:measure}
A compact array of scale $L\ge1$ is a finite sum
\begin{equation}
 U(x,t)=\sum_\alpha c_\alpha L^{-n/2}
 e^{i(x\cdot\omega_\alpha-t|\omega_\alpha|^2)}
 \phi_\alpha\!\left(
     \frac{x-y_\alpha-2t\omega_\alpha}{L},\frac t{L^2}\right).
 \label{d:array}
\end{equation}
The frequencies $\omega_\alpha$ belong to a fixed bounded box.
The native profiles $\phi_\alpha$ are supported in a fixed box and have
uniform bounds for derivatives through a prescribed sufficiently large
finite order.  The number of labels
$(L\omega_\alpha,y_\alpha/L)$ in every unit phase-space ball is
bounded by a fixed constant.  Its coefficient energy is
$E(U)=\sum_\alpha|c_\alpha|^2$.
Subarrays retain their original profiles and labels.

A nonnegative Borel measure $\mu$ has $n$-dimensional growth with
parameter $m$ if
\begin{equation}
 \mu(B(z,r))\le m r^n\qquad(z\in\R^{n+1},\ r\ge1).
 \label{d:ball-growth}
\end{equation}
At array scale $L$ its native time range is bounded by a fixed
constant, so $\supp\mu\subset\{|t|\le C L^2\}$ after a common time
translation.  Fix a sufficiently large $M$.  The tube moment condition
with sparsity parameter $D$ is
\begin{equation}
 \int\left(1+
       \frac{|x-y_\alpha-2\omega_\alpha t|}{L}\right)^{-M}
          \dd\mu(x,t)
       \le \frac{mL^{n+1}}D
       \quad\text{for every included label }\alpha.
 \label{d:tube-moment}
\end{equation}
We use half-open unit cubes for all measure sums, so their masses add
without boundary ambiguity.
\end{definition}

To compare these conventions with \Cref{p:conventions}, set $k=n$
and $N=L$, and let $C_{\rm prof}$ be a fixed common bound for the
required native-profile derivatives. The choice
$s_\alpha=C_{\rm prof}L^{-n/2}|c_\alpha|$ is admissible, so the packet
energy parameter in \eqref{p:energy} is $C_{\rm prof}^2L^2E(U)$.
This is a normalization of packet sizes, not a spacetime
orthogonality assertion. The parameter $D$ in
\eqref{d:tube-moment} measures sparsity.

At any fixed bounded native time, transporting the intercepts along
their velocities preserves phase-space multiplicity: in the normalized
label variables it is a linear shear with bounded operator norm and
bounded inverse.  In particular,
\begin{equation}
 \norm{U}_\infty\le C E(U)^{1/2}.
 \label{d:array-sup}
\end{equation}
Indeed only $O(L^n)$ frequency/intercept labels can meet a point in the
fixed profile support; Cauchy--Schwarz cancels their square-root count
against $L^{-n/2}$.
Condition \eqref{d:tube-moment} with $D=1$ follows from
\eqref{d:ball-growth}, with fixed inflation of $m$.  To see this, on the
annulus of transverse radius $r=2^aL$, the bounded-time tube is covered
by $C(1+L^2/r)$ balls of radius $Cr$.  Its mass is at most
$Cm(r^n+L^2r^{n-1})$.  Multiplication by $2^{-aM}$ and summation over
$a\ge0$ give $CmL^{n+1}$ when $M$ is sufficiently large.

For later use, if labels lie in a frequency cube of side $H^{-1}$
centered at $v$, put
\begin{equation}
 X_{H,v}(x,t)=\left(\frac{x-2vt}{H},\frac t{H^2}\right).
 \label{d:parabolic-map}
\end{equation}
After removal of the common carrier, the field is
$H^{-n/2}U'(X_{H,v}(x,t))$, where $U'$ has scale $L/H$, frequency
$\omega'_\alpha=H(\omega_\alpha-v)$, intercept $y'_\alpha=y_\alpha/H$,
the same $c_\alpha$, and exactly the same native profile.  Thus native
regularity and normalized phase-space multiplicity are unchanged.

\subsection{The precise fractal input and its weighted form}

Write
\[
 \mathcal E g(x,t)=(2\pi)^{-n}
      \int e^{ix\cdot\xi-it|\xi|^2}g(\xi)\dd\xi,
\]
where $g\in L^2$ is supported in a fixed bounded frequency box.
Fixed changes of that box alter constants only.

\begin{theorem}[Du--Zhang input in squared form]
\label{d:dz-input}
Let $\mathcal Q$ be a collection of distinct unit lattice cubes in a
ball of radius $S\ge1$.  Suppose their counting function in every ball
of radius $r\ge1$ is at most $\Gamma r^n$, and their number in any ball
of radius $\sqrt S$ is at most $\lambda_0$.  Then for every $\eps>0$,
\begin{equation}
 \sum_{q\in\mathcal Q}\sup_q|\mathcal E g|^2
 \le C_\eps S^{\eps}
       \Gamma^{4/((n+1)(n+2))}\lambda_0^{\gamma_n}
       S^{\gamma_n}\norm{g}_2^2.
 \label{d:dz-sup}
\end{equation}
All fixed enlargements and translations of the cube lattice are
allowed, with fixed changes of the constant.
\end{theorem}

\begin{proof}
The exact external input is \citet[Theorem~1.6]{DuZhang2019} with counting
exponent $n$.  On a union $X$ of unit cubes for which every occupied
$\sqrt S$ lattice parent has comparable occupancy $\lambda$, it states
\[
 \norm{\mathcal E g}_{L^2(X)}
 \le C_\eps\Gamma^{2/((n+1)(n+2))}
       \lambda^{n/((n+1)(n+2))}
       S^{n/((n+1)(n+2))+\eps}\norm{g}_2.
\]
The original unit-ball frequency hypothesis is changed to our fixed
box by a fixed parabolic dilation and bounded coverings.  The same
holds for balls with arbitrary centers by translating the extension,
which only modulates its input.  Squaring and renaming $\eps$ gives
the exponents in \eqref{d:dz-sup} for an integral over $X$.
No comparable-field-norm assumption occurs in this theorem; that
additional hypothesis belongs to Proposition~3.1 of the cited paper,
which is not needed here.

Choose a comparable ambient scale with a compatible square-root grid.
Group whole square-root parents according to their dyadic occupancy.
A parent is covered by boundedly many radius-$\sqrt S$ balls, so its
occupancy is at most $C\lambda_0$.  Each group inherits the counting
bound.  Sum the squared theorem over these groups; the resulting series
$\sum_{2^a\le C\lambda_0}2^{a\gamma_n}$ is bounded by
$C\lambda_0^{\gamma_n}$.  This removes the equal-occupancy hypothesis.

Finally, $\mathcal E g$ has spacetime Fourier support in a fixed compact
paraboloid patch.  Choose a Schwartz reproducing kernel equal to one
on that patch in Fourier space.  Cauchy--Schwarz and local boundedness
of translates of this kernel imply, for every sufficiently large $A$,
\[
 \sup_q|\mathcal E g|^2
 \le C_A\int_{\R^{n+1}}(1+|h|)^{-A}
            \int_q|\mathcal E g(z+h)|^2\dd z\dd h.
\]
For each fixed $h$, the translated field is an exact extension of an
input of the same norm and support.  Apply the integral estimate on
the unchanged cube collection, then integrate in $h$.  This proves
\eqref{d:dz-sup} without enlarging the ambient radius for distant
translations.
\end{proof}

\begin{proposition}[Weighted fractal norm and occupancy gain]
\label{d:fractal}
\label{d:refined}
Suppose $\mu$ is supported in a ball of radius $CS$, satisfies
\eqref{d:ball-growth}, and, in addition,
\begin{equation}
 \mu(B(z,\sqrt S))\le d\,m S^{n/2},
 \qquad S^{-A_0}\le d\le1.
 \label{d:sqrt-occupancy}
\end{equation}
Then
\begin{equation}
 \sum_q\mu(q)^\beta\sup_q|\mathcal E g|^2
 \le C_{\eps,A_0}\,m^\beta
       S^{n/(n+1)+\eps}d^{\gamma_n}\norm{g}_2^2.
 \label{d:weighted-fractal}
\end{equation}
With $d=1$ no extra assumption beyond growth is needed.  The estimate
is homogeneous in $\mu$ and in the squared input norm.
\end{proposition}

\begin{proof}
Normalize $m=\norm{g}_2=1$; zero cases are immediate.  Sort unit cubes
by $a<\mu(q)\le2a$, where $a$ runs over a fixed multiple of the
nonpositive dyadic powers.  Cubes meeting a radius-$r$ ball lie in a
fixed enlargement, and their disjoint masses are at least $a$ times
their number.  Thus this class has counting parameter at most $Ca^{-1}$
and square-root occupancy at most $CdS^{n/2}a^{-1}$.  Applying
\Cref{d:dz-input} and multiplying by $(2a)^\beta$ gives
\begin{align*}
 \sum_{a<\mu(q)\le2a}\mu(q)^\beta\sup_q|\mathcal E g|^2
 &\le C_\eps S^\eps a^\beta
       (a^{-1})^{4/((n+1)(n+2))}
       (dS^{n/2}a^{-1})^{\gamma_n}S^{\gamma_n}\\
 &=C_\eps S^{n/(n+1)+\eps}d^{\gamma_n}.
\end{align*}
The two cancellations are
\[
 \beta-\frac4{(n+1)(n+2)}-\gamma_n=0,
 \qquad (1+n/2)\gamma_n=\frac n{n+1}.
\]
This exact cancellation leaves no positive power of $a$ to sum, so a
small-mass truncation is necessary.  Bounded frequency support gives
$\norm{\mathcal E g}_\infty\le C$, and there are $O(S^{n+1})$ relevant
unit cubes.  Those with $\mu(q)\le S^{-A}$ contribute at most
$CS^{n+1-A\beta}$.  Choose $A$ so large, in terms of $A_0$ and a
prescribed error power, that this is bounded by
$S^{-B}d^{\gamma_n}$.  There remain $O_A(1+\log S)$ mass classes.
Use half the requested loss in the preceding class estimate and absorb
this count into the remaining $S^\eps$.  Restoring homogeneity proves
\eqref{d:weighted-fractal}.
\end{proof}

\subsection{Exact Fourier layers and spatial localization}

The following reduction permits compact native profiles in estimates
that tolerate an arbitrarily small power loss.  Its common frequency
shift is useful when the chosen subarray changes from one test to
another.

\begin{lemma}[Exact Fourier layers of compact profiles]
\label{d:layers}
For every prescribed finite collection of derivative and native spatial
decay orders, a sufficiently high finite profile regularity in
\eqref{d:array} gives a representation
\begin{equation}
 U(x,t)=\sum_{k\in\Z^n}\int_\R
       \rho(k,u)e^{iu t/L^2}\mathcal E g_{k,u}(x,t)\dd u,
 \label{d:layer-representation}
\end{equation}
where $\rho(k,u)\le C_A(1+|k|+|u|)^{-A}$ for any prescribed $A$.
For every subarray $\mathcal A$,
\begin{equation}
 \norm{g_{k,u}^{\mathcal A}}_2^2
       \le C\sum_{\alpha\in\mathcal A}|c_\alpha|^2.
 \label{d:layer-subarray-energy}
\end{equation}
Each label contributes an exact input supported within $O(L^{-1})$ of
$\omega_\alpha+k/L$, with uniformly bounded native frequency
seminorms.  On a fixed native time range, its extension has all the
prescribed spatial decay about the original trajectory after the
factor $\rho$ is extracted.  The shift $k/L$ and the time modulation
are common to the entire layer and every subarray.  For any fixed
$\eta>0$, the portion $|k|>L^\eta$ can be made an arbitrarily negative
power error in the array supremum, and hence in tests involving
polynomially many cubes and polynomially bounded mass parameters.
\end{lemma}

\begin{proof}
Fourier transform $\phi_\alpha$ in its native variables, denoting the
dual variables by $(\eta,\zeta)$.  Insert a smooth bounded-overlap
partition $\sum_k\vartheta_k(\eta)=1$ into bounded boxes about
$k\in\Z^n$, and change variables $u=\zeta+|\eta|^2$.
For fixed $(k,u)$ define the input, in $\xi=\omega_\alpha+\eta/L$
coordinates, by
\[
 g_{\alpha,k,u}(\omega_\alpha+\eta/L)
   =C c_\alpha L^{n/2}e^{-iy_\alpha\cdot\eta/L}
       \vartheta_k(\eta)
       \widehat\phi_\alpha(\eta,u-|\eta|^2).
\]
Fourier inversion gives
\[
 U=\sum_k\int e^{iu t/L^2}
                     \mathcal E\Big(\sum_\alpha g_{\alpha,k,u}\Big)
                     \dd u.
\]
Compact support and sufficiently many derivatives of $\phi_\alpha$
give any prescribed polynomial decay in $(\eta,\zeta)$, also after a
fixed number of frequency derivatives.  The change
$\zeta=u-|\eta|^2$ costs only polynomial factors.  Extract
$\rho(k,u)=C_A(1+|k|+|u|)^{-A}$ with $C_A$ chosen to bound the needed
seminorms, and put the remaining normalized sum into $g_{k,u}$.
Increasing the finite initial derivative order supplies these bounds
simultaneously.

The input norm estimate follows by a summable Gram-matrix bound.
Two normalized input bumps can interact only if
$|L(\omega_\alpha-\omega_{\alpha'})|\le C$, since the shift $k$ is
common.  For these pairs, integration by parts in the native frequency
variable gives any prescribed decay in
$|y_\alpha-y_{\alpha'}|/L$.  The coefficient factors are
$c_\alpha\overline{c_{\alpha'}}$ and all remaining seminorms are
uniform.  Phase-space counting therefore bounds the absolute row and
column sums of the Gram matrix.  Schur's inequality proves
\eqref{d:layer-subarray-energy}, without requiring the subarray to be a
Fourier projection.

For spatial decay, the phase in native variables is
$z\cdot\eta-\tau|\eta|^2$.  On bounded $\tau$ and $\eta$ near $k$,
integration by parts gives rapid decay in $z$ away from the original
trajectory, at the cost of a fixed polynomial in $k$.  This was
included in the orders used to extract $\rho$; derivatives of the
extension are handled at the same time.  The common shift preserves
all relative frequency labels and phase-space counts.
Finally, Cauchy--Schwarz over labels and the prescribed rapid spatial
decay bound every normalized layer by a fixed polynomial in $(k,u)$
times $E(U)^{1/2}$.  Increase $A$, sum over $|k|>L^\eta$, and integrate
in $u$ to obtain $C_B L^{-B}E(U)^{1/2}$ for any required $B$.
All orders here are fixed before the varying scale $L$.
\end{proof}

\begin{lemma}[Exact packet localization on a time slab]
\label{d:packet-expansion}
Let $T\ge1$ and let $g$ have bounded frequency support.  On a time
interval of length $O(T)$, an exact extension can be localized to
spatial boxes of side $T$ with summed input energies
$C\norm{g}_2^2$.  The localization errors have arbitrarily high negative
powers of $T$, with summable local squared budgets.  No factor equal
to the number of spatial boxes is needed.
\end{lemma}

\begin{proof}
Set $R=\sqrt T$ and absorb the initial time into the input by a
unit-modulus multiplier.  Partition frequency into $R^{-1}$ caps and
expand each rescaled cap density in a Fourier series on a fixed larger
box.  With smooth bumps equal to one on the cap supports, this gives
exact packets
\[
 c_{\theta,m}R^{-n/2}e^{i(x\cdot\omega_\theta-t|\omega_\theta|^2)}
 H_\theta\!\left(
     \frac{x-\gamma Rm-2t\omega_\theta}{R},\frac t{R^2}\right),
 \qquad \sum_{\theta,m}|c_{\theta,m}|^2\le C\norm{g}_2^2.
\]
Here $H_\theta$ is the extension of a fixed smooth native frequency
bump.  It and any fixed derivatives decay arbitrarily rapidly in the
first variable for bounded second variable.  Parseval proves the
coefficient bound; differentiating the input density is unnecessary.
The same cap-overlap and integration-by-parts Gram bound used in
\Cref{d:layers} bounds the squared norm of any retained exact input by
its coefficient energy.

For a spatial box $B$ of side $T$, retain packets whose trajectories
on the time slab come within $CT$ of $B$.  A bounded-velocity
trajectory of duration $O(T)$ can be retained for only $O(1)$ such
boxes.  Therefore these retained energies sum to
$C\norm{g}_2^2$.  For the omitted packets, the native distance from the
box is at least $c\sqrt T$.  Extract an arbitrary negative power of
$T$ from the rapid packet decay and keep a further weight in the
distance from the trajectory to $B$, measured in units of $T$.
Cauchy--Schwarz and phase-space counting bound the squared supremum of
the omitted field on $B$ by
\[
 C_A T^{-A}\sum_{\theta,m}|c_{\theta,m}|^2
       \left(1+\frac{\dist(B,\text{trajectory}_{\theta,m})}{T}
       \right)^{-n-2}.
\]
The right-hand budgets sum boundedly in $B$.  This proves the asserted
error control and also permits arbitrarily long spatial unions of
such boxes.  Translations modulate the exact inputs and preserve their
norms, so the construction is uniform in the slab and box centers.
\end{proof}

\begin{corollary}[Power-loss norm for compact arrays]
\label{d:power-loss}
Let $U$ be a compact array of scale $L$, and let $\mu$ satisfy
\eqref{d:ball-growth} in a spacetime box of side $O(L^2)$.
For every $\eps>0$, with sufficiently many fixed profile derivatives,
\begin{equation}
 \sum_q\mu(q)^\beta\sup_q|U|^2
       \le C_\eps m^\beta L^{2n/(n+1)+\eps}E(U).
 \label{d:compact-fractal}
\end{equation}
\end{corollary}

\begin{proof}
Apply \Cref{d:layers}, discarding $|k|>L^\eta$ with a sufficiently high
negative power.  For $0<\eta<1$, every retained layer has support in a
fixed bounded frequency box.  Apply \Cref{d:fractal} with $S=L^2$ and
$d=1$ to each exact layer.  Minkowski's inequality for the weighted
$\ell^2$ norm of unit-cube suprema, followed by
\eqref{d:layer-subarray-energy}, gives the result because
$\sum_k\int\rho(k,u)\dd u<\infty$.  The common time modulation has
absolute value one.  The discarded error is negligible: each cube
has mass at most $Cm$, there are $O(L^{2(n+1)})$ relevant cubes, and
its supremum has any required negative power of $L$.  Rename the
arbitrarily small scale loss after replacing $S$ by $L^2$.
\end{proof}

\subsection{Coarse decoupling and exact transport powers}

For dyadic $K\ge2$, $0<b<1$, and $P=2/(1-b)$, define
\begin{equation}
 \mathcal T_b(V,\nu)
       =\sum_{Q\in\mathcal Q_{K^2}}
              \nu(Q)^b\norm{V}_{L^P(Q)}^2,
 \label{d:coarse-functional}
\end{equation}
where $\mathcal Q_{K^2}$ is a fixed grid of half-open spacetime cubes
of side $K^2$.

\begin{lemma}[Decoupling inside a strip]
\label{d:strip-decoupling}
Group a field $V=\sum_\theta V_\theta$ by frequency-label bins of side
$K^{-1}$.  Suppose the field supports lie in fixed enlargements of
these bins and within vertical Fourier distance $O(K^{-2})$ of the
paraboloid.  If the labels lie within $O(K^{-1})$ of an affine
$j$-plane, $1\le j\le n$, and $2\le P\le q_j=2(j+2)/j$, then
\begin{equation}
 \norm{V}_{L^P(Q)}^2
 \le C_{\eps,A}K^\eps
       \sum_\theta\sum_{a\in\Z^{n+1}}\braket{a}^{-A}
                 \norm{V_\theta}_{L^P(Q+K^2a)}^2.
 \label{d:local-strip-decoupling}
\end{equation}
The constituent fields may be arbitrary prescribed subfields with
these supports.  Without the strip hypothesis the same estimate holds
with the factor $CK^n$ in place of $C_\eps K^\eps$.
\end{lemma}

\begin{proof}
Rotate frequency into tangent and normal coordinates for the affine
plane, remove its central normal carrier, and shear away the normal
linear time phase.  The remaining normal frequencies have size
$O(K^{-1})$, so their quadratic contribution is $O(K^{-2})$.
On each fixed normal physical slice, the tangential spacetime Fourier
support is therefore within $O(K^{-2})$ of the positive $j$-dimensional
paraboloid.  Multiply first by a Schwartz cutoff bounded below on $Q$
with Fourier width $O(K^{-2})$.  It preserves this support statement
with fixed enlargement.

Apply paraboloid decoupling \citep{BD2015} in the range $P\le q_j$.
After this localization, use smooth tangential cap projections; their
spatial convolution kernels have bounded $L^1$ norm.  Only boundedly
many original enlarged grid bins in the strip can meet a projected
bin, since both the normal thickness and tangential cap size are
$O(K^{-1})$.  Bounded overlap and Minkowski in the normal physical
variables give the squared-norm estimate with a Schwartz weight
outside $Q$.  Decomposing that weight into translates of $Q$ gives
\eqref{d:local-strip-decoupling}, with any fixed $A$.
This proof uses only supports and not the way a subfield was selected.
Without a strip, there are $O(K^n)$ bins, and Cauchy--Schwarz gives the
stated alternative directly.
\end{proof}

\begin{lemma}[Coarse norm transport]
\label{d:transport}
Let $F=X_{K,v}$ with bounded $v$, and suppose
$V(z)=K^{-n/2}e^{i\varphi(z)}V'(Fz)$ with $|e^{i\varphi}|=1$.
For a fixed $a\in\Z^{n+1}$,
\begin{equation}
 \sum_Q\nu(Q)^b\norm{V}_{L^P(Q+K^2a)}^2
 \le C K^{2-(n+2)b}\sum_{r,s\in\mathcal R}
                 \mathcal T_b(V',\nu_{a,r,s}),
 \label{d:transport-functional}
\end{equation}
where $\mathcal R$ is a fixed finite set, independent of $Q,K,a$, and
\begin{equation}
 \nu_{a,r,s}
  =(\tau_{F(K^2a)+K^2(s-r)})_\# F_\#\nu,
 \qquad \tau_h(z)=z+h.
 \label{d:translated-measures}
\end{equation}
Their translations have size at most $CK^2(1+|a|)$.  The same conclusion
holds with any prescribed subset of parent cubes on the left.
If $\nu$ has growth parameter $m$, each child measure has parameter
$CmK^{n+1}$.  Direct pushforward transports
\eqref{d:tube-moment} from scale $L$ to $L/K$ with the same $D$.
\end{lemma}

\begin{proof}
The Jacobian of $F$ is $K^{-(n+2)}$.  Thus, with
$A_Q=F(Q+K^2a)$,
\begin{align*}
 \norm{V}_{L^P(Q+K^2a)}^2
 &=K^{-n+2(n+2)/P}
          \left(\int_{A_Q}|V'|^P\right)^{2/P}\\
 &=K^{2-(n+2)b}
          \left(\int_{A_Q}|V'|^P\right)^{1-b}.
\end{align*}
No density factor accompanies the mass: the measure is the direct
pushforward.  Let $\nu_a=(\tau_{F(K^2a)})_\#F_\#\nu$.
Then $\nu_a(A_Q)=\nu(Q)$, and the $A_Q$ are disjoint.
Group them according to the child $K^2$ cube $C$ containing their
centers.  Each $A_Q$ has diameter $O(K)$, so all sets in a group are
contained in
$C^*=\bigcup_{r\in\mathcal R}(C+K^2r)$ for a fixed finite
$\mathcal R$.  Hölder with exponents $1/b$ and $1/(1-b)$ gives
\[
 \sum_{Q\text{ in the group}}\nu(Q)^b
          \left(\int_{A_Q}|V'|^P\right)^{1-b}
 \le \nu_a(C^*)^b\norm{V'}_{L^P(C^*)}^2.
\]
Subadditivity of the powers $b$ and $1-b$ bounds this by
\[
 \sum_{r,s\in\mathcal R}\nu_a(C+K^2r)^b
                      \norm{V'}_{L^P(C+K^2s)}^2.
\]
The two indices allow the mass and the field norm to occupy different
neighbors.  Reindex by $D=C+K^2s$ to obtain precisely
\eqref{d:translated-measures} and
\eqref{d:transport-functional}.  None of the translations depends on
which parent cube was selected.

The inverse image of a radius-$r$ ball under $F$ has $n$ transverse
widths $O(Kr)$ and longitudinal length $O(K^2r)$.  It is covered by
$O(K)$ balls of radius $CKr$.  Hence its mass is at most
$CmK^{n+1}r^n$ for $r\ge1$, and translations preserve this bound.
For the tube moment, writing $L'=L/K$,
$\omega'=K(\omega-v)$, and $y'=y/K$ gives the identity
\[
 \frac{x'-y'-2\omega't'}{L'}
       =\frac{x-y-2\omega t}{L},
 \qquad (mK^{n+1})(L/K)^{n+1}=mL^{n+1}.
\]
This proves the direct-pushforward assertion.  When the tube centers
are kept fixed, translated measures preserve the moment up to a fixed
factor only if their native displacements are bounded.  Accordingly,
applications requiring this moment truncate the Schwartz shifts while
the child packet widths are much larger than the retained shifts.
The freely chosen exponent $A$ in
\eqref{d:local-strip-decoupling} makes the omitted tail smaller than
any prescribed power once the global suprema and occupied cube counts
are polynomially bounded.  For exact-field applications using only
ball growth, arbitrary translations are harmless.
\end{proof}

\subsection{A power gain from codimension two}

\begin{proposition}[Codimension-two strip saving]
\label{d:strip}
Fix $0<\alpha<1/8$.  Suppose $g$ has bounded frequency support contained
in an $O(S^{-\alpha})$ neighborhood of an affine $(n-2)$-plane, and
$\mathcal Q$ is a collection of unit cubes in a ball of radius $S$ with
at most $m_0r^n$ cubes in every radius-$r$ ball, $r\ge1$.
Fix any $b$ with $\beta<b<2/n$.
There is $c_0>0$, depending only on the fixed dimensional choices and
independent of $\alpha$, such that for every $\eps>0$,
\begin{equation}
 \sum_{q\in\mathcal Q}\sup_q|\mathcal E g|^2
 \le C_{\alpha,\eps}m_0^b
           S^{n/(n+1)-c_0\alpha+\eps}\norm{g}_2^2.
 \label{d:strip-saving}
\end{equation}
In particular, for $m_0=S^{o(1)}$ this is the saving
$S^{n/(n+1)-c_0\alpha+o(1)}$.  Finite disjoint subinputs are allowed,
with their squared norms summed on the right.
\end{proposition}

\begin{proof}
Choose a sufficiently large integer $J$, fixed independently of $S$
and $\alpha$, and dyadic $K$ so that
$H=K^J\asymp S^\alpha$, with fixed rounding factors.
Set $P=2/(1-b)$.  Since $b<2/n$, one has
$P<q_{n-2}=2n/(n-2)$.
Use the counting measure of the selected unit cubes, assigning unit
mass to their centers.  Local constancy and Hölder bound the initial
sum of squared suprema by the coarse functional
\eqref{d:coarse-functional}, with a fixed $K^C$ endpoint cost allowed.
Explicitly, in each coarse cube the sum of the local unit norms is
bounded by the number of selected unit cubes to the power
$1-2/P=b$, times the squared $L^P$ norm on a fixed enlargement.
Schwartz reproduction supplies the required neighboring translates.
Its measure parameter is $Cm_0$.

Iterate \Cref{d:strip-decoupling,d:transport} for $J$ levels.
After $a$ rescalings the strip width is $O(K^a/H)$, which is at most
$O(K^{-1})$ before the last step.  Thus the codimension-two decoupling
hypothesis persists.  The product of transport factors is
$H^{2-(n+2)b}$, up to constants depending on the fixed $J$, while the
final measure parameter is
\[
 m'\le C_Jm_0H^{n+1}.
\]
Smooth frequency decompositions have bounded overlap, so final child
input squared norms sum to $C_J\norm{g}_2^2$.
The $J$ decoupling losses can be made $S^\eps$ by starting with a
sufficiently small loss.  Schwartz shifts may be truncated at a
polynomial radius with negligible error, using bounded-frequency
suprema and polynomial cube counts.  Their weight sums cost fixed
constants at each of the fixed number of levels.

The final time length is
$T=S/H^2$, still a fixed positive power of $S$.
The spatial region can have side $S/H$; to avoid paying for this
aspect ratio, cover it by spatial boxes of side $T$ and use
\Cref{d:packet-expansion} on the child time slab.  Retained exact inputs
have squared norms summable to the child energy.  Their omitted
errors have arbitrarily high negative powers of $T$, with summable
local budgets.  Each resulting spacetime box has side $O(T)$.

Here is the terminal coarse-functional estimate on one such box:
\begin{equation}
 \mathcal T_b(V',\nu')
 \le C_\eps K^{C_*}T^{n/(n+1)+\eps}(m')^b
                      \norm{\text{child input}}_2^2.
 \label{d:terminal-functional}
\end{equation}
Its exponent $C_*$ is independent of $J$.  Indeed, replace the field
norm on each $K^2$ cube by its volume factor times a unit-cube witness
for its supremum.  Move the entire mass of that coarse cube to that
witness.  The resulting measure has ball parameter at most
$CK^{2n}m'$, and each moved mass is at most $CK^{2n}m'$.
Since $b\ge\beta$, its $b$th power is at most
$(CK^{2n}m')^{b-\beta}$ times its $\beta$th power.
Apply \Cref{d:fractal} at scale $T$ to the moved masses.  The volume
factor, the moved growth parameter, and this last comparison produce
only $K^{C_*}$, proving \eqref{d:terminal-functional}.
A bounded number of witness-grid choices and enlargements have the
same cost.  The inequality holds separately for every translated
child measure, and translations are absorbed into modulations of the
exact inputs.  Consequently no count of spatial boxes or translation
labels is introduced.

Combining the three scale contributions, apart from
$S^{n/(n+1)+\eps}$ and the fixed endpoint factor $K^{C_*}$, gives
\begin{align*}
 H^{2-(n+2)b}\,H^{(n+1)b}\,H^{-2n/(n+1)}
   &=H^{2-b-2n/(n+1)}\\
   &=H^{\beta-b}.
\end{align*}
Enlarge $C_*$ to include the initial endpoint cost, still independently
of $J$.  Choose $J$ so large that
$C_*/J<(b-\beta)/2$.  Then
$K^{C_*}H^{\beta-b}\le H^{-(b-\beta)/2}$.
Taking, for example, any $c_0<(b-\beta)/2$ absorbs the fixed rounding
and proves \eqref{d:strip-saving}.  Constants may depend on fixed
$\alpha,J,\eps$, but the saving exponent $c_0$ does not depend on
$\alpha$.  For disjoint subinputs the same construction is applied
with their coefficient energies, whose sum is the original energy.
\end{proof}
